\section{Observational Data}
To test our model against sources where streamers have been detected, we use five observational datasets: Per-emb-2, Per-emb-50, S~CrA, HL~Tau, and SVS~13A. Per-emb-2, Per-emb-50, and SVS~13A were observed with the Northern Extended Millimeter Array (NOEMA), while S~CrA and HL~Tau were observed with the Atacama Large Millimeter/submillimeter Array (ALMA).

Per-emb-2 was observed in D-configuration on June 22, July 27-28, and August 19, 2018, and in C-configuration on October 9, 10, 17, and 18, 2018 (project S18AG; \citealt{Pineda_2020}). The baselines ranged from 16.4 to 287.8~m. Per-emb-50 was observed in C-configuration on December 29, 2019, and January 5, 2021, and in D-configuration on August 6, September 7, and 8, 2020 \citep{Valdivia-Mena_2022}. SVS~13A was also observed in two configurations: C-configuration on December 29, 2019, and January 5, 2020, and D-configuration on August 6 and September 7, 2020, with baselines between 15.6 and 368~m, corresponding to a synthesized beam of $1\farcs2\times0\farcs7$ \citep{Hsieh_2023}.

The S~CrA dataset was obtained from the ALMA Science Archive and was observed on November 6, 2019 (project 2019.1.01792.S; \citealt{Gupta_2023}). HL~Tau was observed with the C40-7 configuration on August 22, 28, and September 1, 2017, with a maximum projected baseline of $\sim$3.3~km \citep{Yen_2019}. For each source, we extract the position-position-velocity (PPV) data of the streamer as described in Section~\ref{sec:fitting_procedure}.

\subsection{Per-emb-2}

\subsection{Per-emb-50}

\subsection{SVS 13A}

\subsection{S CrA}

\subsection{HL Tau}

\section{Build up a model}
We build a model that can reproduce the observed streamer by considering a self-collapse density enhancement structure with a rotation velocity and evolve with time in 3D space. 
\section{Formalism} \label{Sub-sec:Formalism}
In this section, we summarize the formalism of the model, including the density enhancement along the polar angle ($\theta$), the infall velocity ($v_{\rm inf}$), the rotation velocity ($v_{\rm rot}$), the accumulated rotation ($\phi$), and the inclination ($i$).
The sonic length.
\begin{align}
    r_{\rm s} = c_{\rm s} T
\end{align}
Infall velocity
\begin{align}
    v_{\rm inf}=-\sqrt{\frac{2GM_*}{r}}-3.3c_{\rm s} \qquad (c_{\rm s}=200~\mathrm{m\,s^{-1}})
\end{align}
The rotation velocity.
\begin{align}
    \displaystyle v_{\rm rot} = r_{\rm s}\omega\Omega_{\rm kep}\left[\left(\frac{r}{r_{\rm s}}\right)^{-1} + \left(\frac{r}{r_{\rm s}}\right)^{-\frac{1}{3}}\right]
\end{align}
The accumulated rotation.
\begin{align}
    \displaystyle \phi = \phi_0 + t_{\rm s} \omega\Omega_{\rm kep}\sqrt{\frac{2c_{\rm s}^3t_{\rm s}}{GM_*}}\left[\left(\frac{r}{r_{\rm s}}\right)^{-\frac{1}{2}} + \left(\frac{r}{r_{\rm s}}\right)^{-\frac{1}{3}}\right]
\end{align}
\section{Coordinate transformation}
Position
\begin{align}
    \left\{ \begin{array}{ccc}
         x = r\sin{\theta} \cos{\phi}\\
         y = r\sin{\theta} \sin{\phi}\\
         z = r\cos{\theta}\\
    \end{array}    
    \right.
\end{align}
Infall velocity
\begin{align}
    \left\{ \begin{array}{l}
     v_{x_{\rm inf}} = v_{\rm inf}\sin{\theta} \cos{\phi} \\
     v_{y_{\rm inf}} = v_{\rm inf}\sin{\theta} \sin{\phi} \\
     v_{z_{\rm inf}} = v_{\rm inf}\cos{\theta} \\
    \end{array} \right.
\end{align}
Rotation velocity
\begin{align}
    \left\{ \begin{array}{l}
         v_x = -v_{\rm rot}\sin{\theta} \sin{\phi} + v_{x_{\rm inf}}\\
         v_y = v_{\rm rot}\sin{\theta} \cos{\phi} + v_{y_{\rm inf}}\\
         v_z = 0 + v_{z_{\rm inf}}\\
    \end{array}
    \right.
\end{align}
\section{Streamer on the sky in the PPV space}
Position rotates by inclination.
Inclination(i): backward (+), forward (-). 
\begin{align}
    \left\{ \begin{array}{l}
     x' = r\sin{\theta} \cos{\phi} \\
     y' = y\cos{i}+z\sin{i} \\
     z' = -y\sin{i}+z\cos{i}=r(-\sin{\theta} \sin{\phi} \sin{i}+\cos{\theta} \cos{i}) \\
    \end{array} \right.
\end{align}

Velocity rotates by inclination
\begin{align}
    \left\{
    \begin{array}{l}
        v_x' = v_x \\[4pt]
        v_y' = v_y\cos i + v_z\sin i
             = v_{\rm rot}\sin\theta \cos\phi \cos i
             + v_{\rm inf}\bigl(\sin\theta \sin\phi \cos i + \cos\theta \sin i\bigr) \\[4pt]
        v_z' = -v_y \sin i + v_z \cos i
             = - v_{\rm rot}\sin\theta\cos\phi\sin i
             + v_{\rm inf}\bigl(-\sin\theta \sin\phi\sin i + \cos\theta\cos i\bigr)
    \end{array}
    \right.
\end{align}

$\theta$ corresponds to the initial polar angle of the streamer, $\phi_0$ is the initial azimuth angle, $i$ is the inclination angle (zero inclination corresponds to an edge-on disk), $\omega$ is the ratio of Keplerian rotation velocity, and $t_{\rm s}$ is the collapsing time of the particle within the streamer.

\section{Fitting procedure}\label{sec:fitting_procedure}
The fitting procedure described below is applied to each dataset to constrain the model parameters. The results are compared to the observed streamer morphology and kinematics to assess the validity of the model.

\subsection{Position angle}
\YN{This step can be neglected if the disk PA on the sky is known. Otherwise, the previous step needs to be repeated until the optimal is found.}

\subsection{Extracting the position and velocity from observations}
To characterize the observed streamer structure, we extract position and velocity centroids as a function of projected radius from the data cube. 
Starting from the masked streamer cube, we define a sequence of radial bins extending along the approximate direction of the streamer. The radial coordinate ($r$) and polar angle ($\theta$) of each pixel relative to the protostar are computed in the image plane, and an initial reference $\theta(r)$ profile is obtained from a smoothed interpolation of visually identified streamer pixels.

For each radial bin, we compute the centroid position using an intensity-weighted average that favors pixels aligned with the local streamer direction. Specifically, we calculate a directional weight
\[
w_\theta = \exp\left[-\frac{(\theta - \theta_{\rm ref}(r))^2}{2\sigma_\theta(r)^2}\right],
\]
where $\theta_{\rm ref}(r)$ is the interpolated streamer direction and $\sigma_\theta(r)$ is derived from the local spatial dispersion of the preliminary centroids. The position centroid ($\bar{x}$, $\bar{z}$) within each radial bin is then computed from
\[
\bar{x} = \frac{\sum I\,w_\theta\,x}{\sum I\,w_\theta}, \qquad
\bar{z} = \frac{\sum I\,w_\theta\,z}{\sum I\,w_\theta},
\]
where $I$ is the data cube intensity at each pixel.

The velocity centroid for each radial bin is computed in the same region using
\[
\bar{v} = \frac{\sum I\,w_\theta\,v}{\sum I\,w_\theta},
\]
where $v$ is the channel velocity converted to the LSR frame. This procedure yields a set of position-velocity centroids that trace the observed streamer morphology. 
These centroids represent a compact, noise-suppressed summary of the streamer trajectory and are used both in the grid search and in the fast-likelihood MCMC analysis. The number of extracted centroids used for each target is summarized in Table~\ref{table:MCMC_configuration}.

\subsection{Grid search parameter estimation}
We model the streamer kinematics using five parameters ($\theta$, $\phi_0$, $i$, $t_{\rm s}$, $\omega$), corresponding to the initial polar and azimuthal angles of the streamer, the inclination of the system, the characteristic collapse time, and the ratio of rotation velocity. 
For each target, we first perform a coarse five-dimensional grid search to identify reasonable regions of parameter space. The grid covers $\theta \in (0, \pi/2)$, $\phi \in (0, 2\pi)$, and $i \in (-\pi/2, \pi/2)$ with uniform sampling, while the collapse time $t_{\rm s}$ t is the characteristic time after the central mass singularity formation, and $\omega$ is sampled linearly between 0 and 1.

The grid search uses the extracted position-velocity centroids to evaluate the model performance. At each grid point, we compute the misfit between the model streamer and the observed centroids using the error metric
\[
E = \sqrt{\left\langle \Delta x^2 + \Delta z^2 + w_v\,\Delta v^2 \right\rangle},
\]
where $\Delta x$ and $\Delta z$ are the projected position differences (in AU), and $\Delta v$ is the line-of-sight velocity residual (in $\rm km\,s^{-1}$). 
The weight
\[
w_v = \left(\frac{\Delta x_{\rm AU}}{\Delta v_{\rm km\,s^{-1}}}\right)^2
\]
is chosen such that spatial and velocity deviations contribute in comparable physical units.

\subsection{Markov Chain Monte Carlo exploration}
After constructing data-driven prior ranges from the grid search, we perform Markov Chain Monte Carlo (MCMC) exploration of the five parameters. The MCMC sampling is carried out using the \texttt{emcee} package, which implements an affine-invariant ensemble sampler that efficiently explores correlated and multi-modal parameter distributions.

We use the grid results to construct the priors for the subsequent MCMC exploration. Specifically, we select all grid models with $E \le (1 + f)\,E_{\min}$, where $E_{\min}$ is the minimum grid misfit and $f$ is typically 0.05-0.10 (relaxed, if necessary, to include at least several tens of models). For each parameter, we take the 5th-95th percentile range of these "good" models and pad this interval by $\sim30\%$ to define a broad, uniform prior. This procedure yields data-driven prior bounds that are wide enough to avoid artificially constraining the solution but narrow enough to exclude clearly disfavored regions of parameter space.

Each MCMC run begins by initializing an ensemble of walkers within the prior intervals, distributed in a narrow region around the best-performing grid-search solution. For the centroid-based likelihood, which evaluates the misfit using the extracted position-velocity centroids, this "fast-likelihood" stage enables efficient refinement of the parameter space and helps to identify the dominant modes of the posterior distribution. The adopted numbers of walkers and steps for each target are summarized in Table~\ref{table:MCMC_configuration}.

For the full datacube fitting, we evaluate the agreement between the model and the observed position-position-velocity structure using a shell-based distance metric. Instead of computing distances for every voxel, we construct a shell-distance cube around the model streamer and compare it with the observed cube through intensity-weighted shell distances. This shell-based likelihood allows the MCMC to probe the full three-dimensional structure while remaining computationally tractable. The corresponding MCMC configurations (walkers, steps, and whether the shell-based likelihood is used) are also listed in Table~\ref{table:MCMC_configuration}.

Convergence of the chains is evaluated using the autocorrelation time of each parameter, as well as visual inspection of walker trajectories. An initial burn-in phase is discarded, and the remaining samples are used to compute posterior medians and credible intervals as the final estimates of the five parameters.

\begin{table}[h]
\centering
\caption{Overview of centroid extraction and MCMC configurations adopted for each target.}\label{table:MCMC_configuration}
\begin{tabular}{lccccc}
\hline\hline
Target & Method & $N_{\rm cent}$ & $N_{\rm walkers}$ & $N_{\rm steps}$ \\
\hline
\multirow{1}{*}{Per-emb-2}
& Centroid-based & 11 & 20 & $2\times10^{4}$ \\
\hline
\multirow{2}{*}{Per-emb-50}
& Centroid-based & 11 & 20 & $2\times10^{4}$ \\
& Shell-based    & -- & 20 & $8\times10^{3}$ \\
\hline
\multirow{1}{*}{S~CrA}
& Centroid-based & 15 & 20 & $8\times10^{4}$ \\
\hline
\multirow{1}{*}{HL~Tau}
& Centroid-based & 11 & 20 & $8\times10^{4}$ \\
\hline
\end{tabular}
\end{table}

\subsection{Fitting with partially known parameters}


% \YN{Re-include the infall term. Perform full model fitting around the optimal parameters found in the first step. }

% \subsection{First fitting: regression to approximative model}
% \YN{Describe how this is allowed by omitting the infall term.}
% Use X, Z, ${V_y}$ to fit with observation data.
% \begin{align}
%     \left\{ \begin{array}{l}
%         X = r\sin{\theta} \cos{\phi} \\
%         Z = r(-\sin{\theta} \sin{\phi} \sin{i}+\cos{\theta} \cos{i})\\
%         V_y = r\Omega_{\rm ref}\left(\frac{r}{r_{\rm ref}}\right)^{-\frac{4}{3}}\sin{\theta} \cos{\phi} \cos{i}+v_{\rm inf}(\sin{\theta} \sin{\phi} \cos{i}+\cos{\theta} \sin{i})
%     \end{array}
%     \right.    
% \end{align}
% Fitting for the initial parameter.
% \begin{align}
%     \left\ \begin{array}{l}
%         \displaystyle\frac{V_y}{X} = \frac{r\sin\theta\Omega_{\rm ref}(\frac{r}{r_{\rm ref}})^{-\frac{4}{3}}\cos\phi\cos i}{r\sin\theta\cos\phi}=\Omega_{\rm ref}(\frac{r}{r_{\rm ref}})^{-\frac{4}{3}}\cos i \\
%         \displaystyle\left(\frac{r}{r_{\rm ref}}\right)^{-\frac{4}{3}}=\frac{V_y}{X\Omega_{\rm ref}\cos i}\\
%         \displaystyle\frac{X}{r}=(\Omega_{\rm ref}\cos i)^{-\frac{3}{4}}\frac{X}{r_{\rm ref}}\left(\frac{V_y}{X}\right)^{\frac{3}{4}}
%     \end{array}
%     \right\
% \end{align}
% \begin{align}
%     \left\ \begin{array}{l}
%         \displaystyle X\cdot\left(\frac{V_y}{X}\right)^{\frac{3}{4}}=r_{\rm ref}\sin\theta\cos\phi (\Omega_{\rm ref}\cos i)^{\frac{3}{4}}\\
%         \displaystyle Z\cdot\left(\frac{V_y}{X}\right)^{\frac{3}{4}}=(\cos\theta\cos i-\sin\theta\sin\phi\sin i)r_{\rm ref}(\Omega_{\rm ref}\cos i)^{\frac{3}{4}}
%         \end{array}
%     \right\
% \end{align}
% \begin{align}
%     \left\ \begin{array}{l}
%         A = r_{\rm ref}(\Omega_{\rm ref}\cos i)^{\frac{3}{4}}\sin\theta \\
%         B =  - r_{\rm ref}(\Omega_{\rm ref}\cos i)^{\frac{3}{4}}\sin\theta\sin i\\
%         C = r_{\rm ref}(\Omega_{\rm ref}\cos i)^{\frac{3}{4}} \cos\theta\cos i \\
%         k = \frac{t}{\sin i} \\
%         \phi = k(R_i + R_0)
%     \end{array}
%     \right\
% \end{align}

% \begin{align}
%     \left\ \begin{array}{l}
%         R_i = \frac{V_y}{X}\\
%         X_i = A\cos(k(R_i+R_0))\\
%         Y_i = B\sin(k(R_i+R_0)) + C
%     \end{array}
%     \right\
% \end{align}

% \begin{equation}
%     \left\ 
%         SEE^2 = \sum(X_i-A\cos{(k(R_i+R_0))})+\sum(Y_i-B\sin(k(R_i+R_0))-C)
%     \right\
% \end{equation}

% \begin{equation}
%     \left\ \begin{array}{l}
%         \displaystyle\frac{1}{2}\frac{\partial SEE^2}{\partial A} = A\sum\cos^2{(k(R_i+R_0))}-\sum X_i\cos(k(R_i+R_0))\\
%         \\
%         \displaystyle\frac{1}{2}\frac{\partial SEE^2}{\partial B} = B\sum\sin^2(k(R_i+R_0)) - \sum Y_i\sin(k(R_i+R_0)) + C\sum\sin(k(R_i+R_0))\\
%         \\
%         \displaystyle\frac{1}{2}\frac{\partial SEE^2}{\partial C} = \sum Y_i - \sum B\sin(k(R_i+R_0))-NC\\
%         \\
%         \begin{split}
%             \frac{1}{2}\frac{\partial SEE^2}{\partial k}=A\sum X_i R_i\sin(k(R_i+R_0)) - A^2\sum R_i\cos(k(R_i+R_0))\sin(k(R_i+R_0))\\
%                                                         - B\sum Y_iR_i\cos(k(R_i+R_0))+B^2\sum R_i\sin(k(R_i+R_0))\cos(k(R_i+R_0))\\
%                                                         + BC\sum R_i\cos(k(R_i+R_0))\\
%                                                         + R_0A\sum X_i \sin(k(R_i+R_0)) - R_0A^2\sum \cos(k(R_i+R_0))\sin(k(R_i+R_0))\\
%                                                         - R_0B\sum Y_i\cos(k(R_i+R_0))+R_0B^2\sum \sin(k(R_i+R_0))\cos(k(R_i+R_0))\\
%                                                         + R_0BC\sum \cos(k(R_i+R_0))\\
%                                                         = A\sum X_i R_i\sin(k(R_i+R_0)) - A^2\sum R_i\cos(k(R_i+R_0))\sin(k(R_i+R_0))\\
%                                                         - B\sum Y_iR_i\cos(k(R_i+R_0))+B^2\sum R_i\sin(k(R_i+R_0))\cos(k(R_i+R_0))\\
%                                                         + BC\sum R_i\cos(k(R_i+R_0))\\            
%         \end{split}
%         \\
%         \begin{split}
%             \frac{1}{2}\frac{\partial SEE^2}{\partial R_0} = kA\sum X_i \sin(k(R_i+R_0)) - kA^2\sum \cos(k(R_i+R_0))\sin(k(R_i+R_0))\\
%                                                             - kB\sum Y_i\cos(k(R_i+R_0)) + kB^2\sum \sin(k(R_i+R_0))\cos(k(R_i+R_0))\\
%                                                             + kBC\sum \cos(k(R_i+R_0))           
%         \end{split}
%     \end{array}
%     \right\
% \end{equation}

% \begin{align}
%     \left\ \begin{array}{l}
%         \displaystyle A = \frac{\sum X_i\cos(k(R_i+R_0))}{\sum\cos^2{(k(R_i+R_0))}}\\
%         \displaystyle B = \frac{N\sum Y_i\sin(k(R_i+R_0)) - \sum Y_i\sum\sin(k(R_i+R_0))}{N\sum \sin^2(k(R_i+R_0)) - \sum\sin(k(R_i+R_0))\sum\sin(k(R_i+R_0))}\\
%         \displaystyle C = \frac{\sum Y_i\sum\sin^2(k(R_i+R_0)) - \sum Y_i\sin(k(R_i+R_0))\sum\sin(k(R_i+R_0))}{N\sum\sin^2(k(R_i+R_0)) - \sum\sin(k(R_i+R_0))\sum\sin(k(R_i+R_0))}
%     \end{array}
%     \right\
% \end{align}

% \begin{align}
%     \left\ \begin{array}{l}
%         \displaystyle i = \arcsin{\frac{-B}{A}}\\
%         \phi_0 = R_0 \cdot k\\
%         \displaystyle \theta = \arctan\frac{A\cos i}{C}\\
%         \displaystyle \Omega_{\rm ref} = \left(\frac{A}{r_{\rm ref}\sin\theta}\right)^{\frac{4}{3}}\frac{1}{\cos i}\\
%         \displaystyle T = k \frac{\cos^{\frac{1}{4}}i}{\Omega_{\rm ref}^{\frac{3}{4}}}
%     \end{array}
%     \right\
% \end{align}

% \YN{The external part of the streamer where the infall term cannot be safely neglected should be truncated. The procedure should be repeated until convergence is obtained.}

% \YN{If the inclination is known, fitting could be simplified. $B=-A\sin{i}$. Re-write the equations for the regression model: only 4 equations to be solved.}